\section{Methods: Group 2 --- Intervertebral Disc}
\label{sec:g2}

Group 2 measures, per disc C2--C3 through C7--T1: disc signal-intensity height, the disc-height index
(DHI), posterior disc bulge, and a cervical signal-degeneration grade. The components form a chain
(\texttt{disc\_si\_height} $\rightarrow$ \texttt{disc\_height\_index} $\rightarrow$ \texttt{disc\_ap\_bulge}
$\rightarrow$ \texttt{pfirrmann\_grade}), threading prior results.

\paragraph{Disc height and DHI.} Disc height is measured at the midline band selected via the canal;
DHI normalizes the disc middle height by the mean of the adjacent vertebral-body middle heights. The
in-code absolute cut (DHI $<0.30$) is uncited (an animal/lumbar borrow); the cervical-appropriate
notion of reduced disc height is a $>$30\% inter-level drop~\cite{suzuki2018} or an absolute height
$<\SI{3}{\milli\meter}$, both requiring cross-level context.

\paragraph{Posterior bulge.} Bulge is the posterior excursion of the disc beyond the vertebral margin.
The correct reference is the tilted chord between adjacent posterior vertebral-body corners; a flat
vertical back-wall mis-measures on lordotic necks. Cervical bulge thresholds derive from
Nakashima~et~al.~\cite{nakashima2015} ($>\SI{1}{\milli\meter}$ bulge); the often-cited
$\SI{1.35}{\milli\meter}$ ``cord-risk'' figure is treated as unverified pending full-text confirmation.

\paragraph{Signal grade.} Cervical disc degeneration is graded by the modified (Miyazaki)
scheme~\cite{miyazaki2008}, computed from a CSF-normalized nucleus signal ratio (requiring the raw T2
volume), \emph{not} the lumbar Pfirrmann scheme. The ratio-to-grade cut-points have no published
$\kappa$/AUC for cervical and are therefore scoped as a research-grade heuristic, surfaced for review.

\paragraph{Validation outcome and a detected bug.} On the cohort, both DHI and bulge read
\bad{backwards} (healthy scored \emph{worse} than the symptomatic cohort; Section~\ref{sec:results}). The
DHI denominator is over-measured at the cervicothoracic junction (adjacent-VB middle heights of
20--27\,mm versus the true $\sim$12--13\,mm), collapsing DHI and over-firing the reduced-disc flag --- the
exact denominator discrepancy predicted by the component's author. Because Group~2 is a teammate-owned
module tuned on a different (Duke) cohort, the bug is documented and flagged with its root cause and a
candidate fix (robustify the adjacent-VB height at junction levels), rather than rewritten without the
author and without ground truth. Group~2 is therefore \emph{not validated} in this pass; the validation
harness correctly identified it as broken.
